% Part: normal-modal-logic
% Chapter: completeness
% Section: introduction

\documentclass[../../../include/open-logic-section]{subfiles}

\begin{document}

\olfileid{nml}{com}{int}

\olsection{प्रस्तावना}

$\Sigma$ ही मोडल प्रणाली असेल, तर निर्दोषता
प्रमेय असे सांगते: $\Sigma \Proves !A$ असल्यास,
ज्या प्रतिमान-वर्ग~$\mClass{C}$ मध्ये~$\Sigma$
मधील सर्व !!{formula}s चे सर्व दृष्टांत वैध
आहेत, त्या वर्गात $!A$ वैध असते.
विशेषतः, $\Log{K}
\Proves !A$ असेल, तर $!A$ सर्व प्रतिमानांत
सत्य असते; $\Log{KT} \Proves !A$ असेल,
तर $!A$ सर्व परावर्ती प्रतिमानांत सत्य असते;
$\Log{KD} \Proves !A$ असेल, तर $!A$
सर्व सर्वत्र उत्तरवर्ती प्रतिमानांत सत्य असते; इत्यादी.

संपूर्णता हा निर्दोषतेचा व्यत्यास आहे:
उदाहरणार्थ, \Log{K} संपूर्ण आहे याचा अर्थ
!!a{formula}~$!A$ वैध असेल, तर $\Proves !A$.
संपूर्णता सिद्ध करणे निर्दोषता सिद्ध करण्यापेक्षा
खूप कठीण आहे. सुरुवातीला प्रतिपरिवर्तन विचारात
घेणे उपयुक्त ठरते: $\Proves/ !A$ असेल तेव्हा
नेहमी प्रतिप्रतिमान, म्हणजे
$\mSat/{M}{!A}$ असे~$\mModel{M}$ प्रतिमान,
अस्तित्वात असेल, तर आणि तेव्हाच \Log{K}
संपूर्ण असते. समतुल्यपणे ($!A$ चा निषेध
घेतल्यास), $\Proves/ \lnot !A$ असेल तेव्हा
$!A$ चे प्रतिमान असते, हे सिद्ध करता येईल.
असे प्रतिमान रचताना $!A$ मधील माहिती
वापरता येते. विशिष्ट !!{formula}s साठी
प्रतिमाने शोधतानाही आपण अनेकदा असेच करतो.
उदाहरणार्थ, $p \lif \Box q$ चे प्रतिप्रतिमान
शोधायचे असेल, तर त्यात $p$ सत्य आणि
$\Box q$ असत्य असलेले जग हवे, हे माहीत असते.
ज्या जगात $\Box q$ असत्य आहे, त्या जगापासून
प्राप्य असे एखादे जग असले पाहिजे की त्यात $q$
असत्य आहे. आपल्याला एवढेच ठरवायचे असते:
कोणत्या जगांत !!{propositional variable}s
सत्य आहेत आणि कोणती जगे कोणत्या जगांपासून
प्राप्य आहेत.

पण संपूर्णता सिद्ध करताना प्रतिमान रचण्यासाठी
आपल्यापाशी एखादे ठरावीक !!{formula}~$!A$
नसते. $\Proves/[\Sigma] \lnot !A$ असलेल्या
प्रत्येक $!A$ साठी प्रतिमान अस्तित्वात आहे,
हे दाखवायचे असते. ही किमान आवश्यक अट आहे:
\emph{जर} $\Proves[\Sigma] \lnot !A$ असेल,
तर निर्दोषतेनुसार~$!A$ चे
($\Sigma$ सत्य असलेले) प्रतिमान असू शकत नाही.
आता लक्षात घ्या की $\Proves/[\Sigma]
\lnot !A$ असेल तेव्हा आणि तेव्हाच $!A$ हे
$\Sigma$-सुसंगत असते. ($\Sigma
\Proves/[\Sigma] \lnot !A$ आणि
$!A \Proves/[\Sigma] \lfalse$ समतुल्य
आहेत, हे आठवा.) म्हणून प्रत्येक
$\Sigma$-सुसंगत !!{formula} साठी प्रतिमान
रचणे हे आपले काम आहे.

आपण पुढील युक्ती वापरू: $!A$ सामावणारा
$\Sigma$-सुसंगत !!{formula}s चा संच शोधू;
त्यात $!A$ सत्य करणारे जग कसे असावे
हे सांगणारी इतर सूत्रेही असतील. असे संच
\emph{संपूर्ण} $\Sigma$-सुसंगत संच असतात.
$!A$ सत्य करण्यासाठी एका जगाचे प्रतिमान
रचणे पुरेसे नाही; त्यात अनेक जगे आणि
प्राप्यता संबंध असावा लागेल. $!A$ सामावणाऱ्या
संपूर्ण $\Sigma$-सुसंगत संचात
$\Box !B$ आणि~$\Diamond !C$ या रूपांतील
इतर !!{formula}s ही असतील. सर्व प्राप्य
जगांत $!B$ सत्य असले पाहिजे; किमान एका
प्राप्य जगात $!C$ सत्य असले पाहिजे. यासाठी
शक्य असलेले \emph{सर्व} संपूर्ण
$\Sigma$-सुसंगत संच जगांच्या संचाचा
आधार म्हणून घेऊ. आपल्या प्रतिमानात एक
संपूर्ण $\Sigma$-सुसंगत संच दुसऱ्यापासून
केव्हा प्राप्य मानायचा, हे ठरवणे अवघड
ठरेल.

असे परिभाषित केलेल्या प्रतिमानात एखाद्या
जगात---जे स्वतः संपूर्ण $\Sigma$-सुसंगत
संच आहे---$!A$ सत्य असेल तेव्हा आणि
तेव्हाच $!A$ त्या संचाचा !!a{element}
असेल, हे दाखवू. $!A$ हे
$\Sigma$-सुसंगत असेल, तर किमान एका
संपूर्ण $\Sigma$-सुसंगत संचाचा ते
!!a{element} असेल (हे आपण सिद्ध करू);
म्हणून $!A$~सत्य असलेले जग मिळेल. अशा
रीतीने प्रत्येक $\Sigma$-सुसंगत
!!{formula}~$!A$ काहीतरी जगात सत्य असलेले
एकच प्रतिमान आपल्याला मिळेल.
हेच~$\Sigma$ साठीचे \emph{कॅनॉनिकल}
प्रतिमान आहे.

\end{document}
